\section{Introduction}\label{sec:intro}

A fault tolerant quantum computer protects its logical qubits with an error correcting code and must decode the
syndrome of that code while the computation runs \cite{dennis2002topological,fowler2012surface}. The rotated surface
code \cite{kitaev2003anyons,fowler2012surface} is the leading candidate for superconducting hardware, where logical memories below threshold have now been
demonstrated \cite{google2023scaling,google2025below}. Its standard decoder, minimum weight perfect matching, is
implemented efficiently on CPUs by PyMatching and the Sparse Blossom algorithm \cite{higgott2022pymatching,
higgott2025sparse}, and union find decoders trade some accuracy for a simpler and faster procedure
\cite{delfosse2021unionfind}. Because a superconducting device produces a new round of syndrome data about every
microsecond, a decoder that falls behind accumulates a backlog that eventually stalls the computation, and dedicated
hardware such as FPGAs has been used to keep up \cite{das2022lilliput,liyanage2023fpga,liyanage2024distributed,
barber2025decoder}. Where syndrome cycles last milliseconds, as in trapped ion architectures, a single CPU can decode
all logical qubits of a large application in real time \cite{ye2026megaquop}; the case for accelerators therefore
rests on the microsecond budget of fast platforms.

Neural networks offer a second route. They have been trained as complete decoders \cite{torlai2017neural,
chamberland2018deep,gicev2023scalable,bausch2024alphaqubit} and, more recently, as \emph{pre-decoders}: a local
network removes the short error chains that dominate the syndrome and leaves a sparser residual syndrome to a global
decoder. Chamberland et al.\ released a family of such pre-decoders, the NVIDIA Ising pre-decoders, as fully
convolutional three dimensional networks over space and time, together with their training code and pretrained
weights \cite{chamberland2026predecoder,nvidia_ising_code,nvidia_fast_card,nvidia_accurate_card}. On GPUs the
pre-decoder both lowers the logical error rate of uncorrelated matching and shortens the time PyMatching needs for the
residual syndrome \cite{chamberland2026predecoder}. A convolutional network of fixed size, applied to a fixed window of
syndrome data, is also a natural workload for an FPGA: its arithmetic is regular, its weights fit on chip, and
integer arithmetic suffices for inference \cite{jacob2018quantization,zhang2015optimizing,umuroglu2017finn,
duarte2018hls4ml}. Overwater et al.\ analysed the hardware cost of neural decoders for this reason
\cite{overwater2022neural}.

This paper places the two published surface code pre-decoders on four AMD Alveo U55C cards. Following our earlier
multi FPGA work, each card is an independent compute node driven by its own host process \cite{ali2026multifpga}; a
batch of syndrome windows is spread over the nodes, which is the shot parallel mode of a decoder serving many
independent windows. The contributions are:

\begin{itemize}
\item An integer model of both networks (8 bit weights and activations, 32 bit accumulation, integer requantization
and a 256 entry GELU table) that keeps NVIDIA's evaluation pipeline unchanged, so that the floating point networks,
the integer model and the FPGA output are judged by the same logical error rate code.
\item One HLS kernel that runs the four layer Fast network and the six layer Accurate network from the same bitstream,
with all activations and partial sums on chip, and that is bit exact with the integer model in C simulation and on
every instance of every card.
\item Measurements on four cards: throughput on one, two and four nodes and latency per window, against the same
networks in PyTorch on a 256 thread host CPU, and logical error rates of the floating point and integer networks
against PyMatching alone.
\item A reproduction of NVIDIA's published accuracy for both networks with their own evaluation code, which shows when a
pre-decoder pays off: the small Fast network gives little or no gain at a physical error rate of $p=0.003$, and the
gains reported for the family come from the larger network and from higher noise.
\end{itemize}

Section~\ref{sec:related} reviews decoders on FPGAs and neural decoders, Section~\ref{sec:background} the
pre-decoder and its evaluation, Section~\ref{sec:quant} the integer model, Section~\ref{sec:arch} the kernel and the
node runtime, Section~\ref{sec:setup} the experimental setup, Section~\ref{sec:results} the results and
Section~\ref{sec:discussion} their limits.
